\chapter{Example sessions} \label{app:examples}

%% DRAFT (Claude, 2026-10-28), rewrite in your own words.
%% Gold traces: data/eval/results/current-20261020T044840.json (frozen code
%% 0293a0b), analysis timings from dense_only-20261020T022900.json. The q-21
%% omission: the workplace creche line sits under ###### NATALIDADE in
%% data/1_markdown/ad/ad-2024-manual.md (retrieved, rank 4); the higher-education
%% creche line under ###### ESTUDANTES: APOIOS SOCIAIS E BEM-ESTAR (not retrieved).
%% Transcripts: copies of data/eval/extras/walkthrough/*.txt, kept here so the
%% Overleaf zip carries them.
This appendix completes Section~\ref{sec:results-practice}. The first part
traces two more gold questions from the test run, as q-03 was traced in
Section~\ref{subsec:results-trace}. The second gives the seven command-line
transcripts of Table~\ref{tab:walkthrough} in full, as the system printed
them.

\section{Two more gold questions} \label{sec:app-traces}

\paragraph{A comparison (q-21).} ``Em que diferem AD e PS nas
creches?''\footnote{``How do the AD and the PS differ on nurseries?''} The
analyser read it as a comparison of the AD and the PS on \textbf{educacao}.
The gold set also expects \textbf{seguranca\_social}, and it marks the
question as asking for the best match rather than everything each party says,
so topic F1 is 0.667 and coverage accuracy 0. The search fetched 8 chunks
per party, 16 in all, merged into 13 passages.
Three of the five gold chunks were among them, at ranks 1, 2 and 4 (recall
0.60, nDCG 0.699). The answer follows in full.

\begin{Verbatim}[fontsize=\scriptsize, breaklines=true, frame=single]
### Aliança Democrática

A Aliança Democrática propõe integrar a faixa etária dos 0 aos 3 anos no sistema educativo tutelado pelo Ministério da Educação [1]. O partido defende o acesso universal e gratuito às creches e ao pré-escolar, pretendendo alargar a oferta pública e sem custos para as famílias através do aumento da capacidade do Estado ou da contratualização com os setores social, particular e cooperativo [1][2]. No plano orçamental, o partido prevê a inclusão de despesas adicionais com as creches e o pré-escolar no seu cenário de despesa [6].

### Partido Socialista

O Partido Socialista compromete-se a prosseguir a política de alargamento da gratuitidade das creches e do pré-escolar [9]. O partido pretende reforçar a capacidade instalada da rede de respostas de creche em todo o país para garantir o acesso gratuito a todas as crianças, visando permitir a conciliação entre a vida pessoal, familiar e profissional [9]. O programa menciona que a gratuitidade das creches é uma medida já em implementação [13].

### Síntese

Ambos os partidos defendem o acesso universal e gratuito às creches e ao pré-escolar [1][2][9][13]. A diferença reside na abordagem de implementação: a Aliança Democrática foca na integração da faixa dos 0-3 anos no sistema educativo tutelado pelo Ministério da Educação e na mobilização dos setores público, social e privado [1][2], enquanto o Partido Socialista enfatiza a continuidade de uma política de alargamento já em curso, visando reforçar a capacidade instalada da rede [9][13].
\end{Verbatim}

Every claim in the answer is in the excerpts (faithfulness 1.00), and it
covers 9 of the reference's 11 claims (correctness recall 0.818). The two it
misses fail for different reasons. Crèches in higher education institutions
are proposed in the AD's section on student support, which the search never
reached, so the generator could not know about them. Tax benefits for firms
that open workplace crèches are proposed in the AD's section on the birth
rate, which the search did retrieve, and the generator left them out. One miss
belongs to retrieval and the other to generation, and only reading the
excerpts tells them apart.

\paragraph{A party outside the corpus (q-59).} ``O que propõe o ADN para a
habitação?''\footnote{``What does the ADN propose for housing?''} The ADN
ran in 2024 but its manifesto is not in the corpus. The analyser recognised
the party, found it outside the corpus and named no corpus party, so
\sysname{} skipped the search and the generator and returned its fixed
answer (Section~\ref{subsec:queryanalysis}).

\begin{Verbatim}[fontsize=\scriptsize, breaklines=true, frame=single]
O corpus não inclui o programa eleitoral destes partidos: ADN. Inclui apenas os de: Aliança Democrática, Bloco de Esquerda, CHEGA, Iniciativa Liberal, LIVRE, Pessoas-Animais-Natureza, Partido Comunista Português, Partido Socialista.
\end{Verbatim}

The answer\footnote{``The corpus does not include the manifesto of these
parties: ADN. It only includes those of'' the eight parties listed.} took one
LLM call, the analysis, and about a second. Before this guard existed,
\sysname{} answered the question with the AD's housing proposals.

\section{Command-line transcripts} \label{sec:app-transcripts}

Each transcript is the full output of one \texttt{python main.py query} call
on the frozen code, unedited. The lines in brackets are the progress messages
shown while the user waits, the answer sits between the two rules, and the
sources and model calls follow it. Each call loads the search model, because
each question ran in a fresh process.

\subsection*{O que propõe o LIVRE para a habitação?}
\VerbatimInput[fontsize=\scriptsize, breaklines=true, frame=single]{appendices/walkthrough/01-single-party.txt}

\subsection*{Compara as propostas do BE e da IL para o Serviço Nacional de Saúde.}
\VerbatimInput[fontsize=\scriptsize, breaklines=true, frame=single]{appendices/walkthrough/02-comparison.txt}

\subsection*{Que partidos falam sobre o bem-estar animal?}
\VerbatimInput[fontsize=\scriptsize, breaklines=true, frame=single]{appendices/walkthrough/03-topic.txt}

\subsection*{Que propostas existem para combater a corrupção?}
\VerbatimInput[fontsize=\scriptsize, breaklines=true, frame=single]{appendices/walkthrough/04-general.txt}

\subsection*{O que propõe o PCP sobre criptomoedas?}
\VerbatimInput[fontsize=\scriptsize, breaklines=true, frame=single]{appendices/walkthrough/05-refusal.txt}

\subsection*{Quais são os planos da AD para a habitação?}
\VerbatimInput[fontsize=\scriptsize, breaklines=true, frame=single]{appendices/walkthrough/06-followup-1.txt}

\subsection*{E o PS?}
\VerbatimInput[fontsize=\scriptsize, breaklines=true, frame=single]{appendices/walkthrough/07-followup-2.txt}
